% =====================================================================================
% Chapter 3 · Uplift targeting — who to treat
%
% Built from notebooks/01_uplift_targeting.ipynb. NOT ONE NUMBER IN THIS FILE IS TYPED BY
% HAND: every figure in the prose is a macro from build/macros.tex, generated by cmp.report
% from the executed notebook. The only literals below are (i) the constants of the
% data-generating model in §3.2, which are *definitions* rather than results, and (ii) the
% decision-rule thresholds (0.8, 0.9), which are *conventions* stated in the text.
%
% A STANDING WARNING TO ANY FUTURE EDITOR. pymc-bart is not seed-reproducible in this
% environment: two consecutive FULL runs of identical source have moved BCF's PEHE by more
% than a fifth and flipped a coverage winner line. So NO BART-derived quantity — PEHE, AUUC,
% coverage, correlation — may ever be written as a literal here, and no VERDICT about one may
% be argued in prose. The notebook computes each head-to-head against twice its own paired
% bootstrap standard error and emits the winner WORD, which is "too close to call" whenever
% that is what the measurement says. This chapter cites those words. Do not replace them with
% a sentence of your own; the next re-execution will make it a lie.
% =====================================================================================
\chapter{Uplift targeting: who to treat}
\label{ch:up}

\begin{quote}\small
\emph{A discount email costs \eur{\nbOneCost} to send to one customer, and the list holds
\nbOneNCustomers{} customers. Its average effect is \eur{\nbOneTrueAte} per customer --- positive, and
recovered to within a euro by three classical estimators (\eur{\nbOneOlsEst}, \eur{\nbOneLinEst},
\eur{\nbOneAipwEst}). That average can issue only two instructions, and both are wrong. The naive
comparison reads \eur{\nbOneNaiveEst}, sits above the cost, and says \emph{email everyone}; the campaign
then \textbf{destroys \eur{\nbOneTreatAllLoss}}. The correct average sits below the cost and says
\emph{email nobody}, which earns nothing and forgoes the whole \eur{\nbOneOracleProfit} on the table.
Estimating the average better fixes neither instruction. \pc{\nbOneShareProfitable} of the list clears
the cost of a send; \pc{\nbOneShareSleeping} are \emph{sleeping dogs} the email actively harms. The
business needs the \emph{ranking} the average hides, not the average. This chapter builds that ranking,
grades it, and prices it: a per-customer rule that captures \pc{\nbOneCapture} of the
\eur{\nbOneOracleProfit} on offer, on a decision where emailing everyone instead \textbf{loses
\eur{\nbOneTreatAllLoss}}.}
\end{quote}

% =====================================================================================
\section{The decision}
\label{sec:up:decision}

A retailer holds a list of \nbOneNCustomers{} customers and a discount offer. Reaching one customer
costs \eur{\nbOneCost}, all in --- the creative, the send, and the discounted margin. The question is
not whether to run the campaign. It is \textbf{whom to send it to}.

The reflex is to estimate the average effect of the email on spend and compare it with the cost.
\Cref{sec:up:classical} does this in full. The average effect is \eur{\nbOneTrueAte}, well above zero
and well below the \eur{\nbOneCost} cost of a send. As a guide to action it is nearly worthless.
Treating everyone earns $n\,(\text{ATE} - c)$, so the break-even \emph{is} the average itself. Below
\eur{\nbOneTrueAte} a contact it says send to everyone; above it, send to nobody. Here it says send to
nobody, which earns zero. Move the cost a little and it flips to \emph{send to everyone}, which
\cref{tab:up:rules} prices at a loss of \eur{\nbOneTreatAllLoss}. Both instructions are wrong. Neither
is repaired by a more precise average: \pc{\nbOneShareProfitable} of the list clears the
\eur{\nbOneCost} cost while \pc{\nbOneShareSleeping} are actively harmed, and one number cannot serve
two groups that want opposite treatment.

Customers are not alike, and the way they differ is the whole commercial opportunity, not a nuisance
to average over. Sort them by what they would do in each future --- emailed and not emailed --- and
four types appear.

\begin{center}\small
\begin{tabular}{@{}lll@{}}
\toprule
 & \textbf{buys if emailed} & \textbf{does not buy if emailed} \\
\midrule
\textbf{buys if not emailed}
  & \emph{Sure Thing} ($\tau \approx 0$)
  & \emph{Sleeping Dog} ($\tau < 0$) \\
\textbf{does not buy if not emailed}
  & \emph{Persuadable} ($\tau > 0$)
  & \emph{Lost Cause} ($\tau \approx 0$) \\
\bottomrule
\end{tabular}
\end{center}

Only the \emph{persuadables} are worth paying to reach: the email tips them over. The \emph{sure
things} would have bought anyway, so the discount is margin given away. The \emph{lost causes} were
never going to buy, so the send is wasted. The \emph{sleeping dogs} make this chapter necessary: for
them the effect is \textbf{negative}. Remind a disengaged customer of a subscription they had stopped
thinking about, and you prompt an unsubscribe, a return, a cancelled auto-renewal. Emailing a sleeping
dog is worse than wasteful. It \emph{destroys} value that would otherwise exist.

Three consequences follow, and they organise the chapter.

\textbf{First, a response model is the wrong tool.} A model that predicts \emph{who buys after being
emailed} cannot tell a persuadable from a sure thing: both bought. The two differ only in their
\emph{effect}, and the effect is never observed.

\textbf{Second, the average is the wrong estimand.} A positive average is fully compatible with a
segment where the campaign destroys money, and here it is: \pc{\nbOneShareSleeping} of the list has
$\tau_i < 0$, and their planted effects sum to \nbOneSleepingLoss{} euros the campaign would erase if
it reached them. The \eur{\nbOneTrueAte} average is real, correctly estimated, and hides them. This is
why \Cref{ch:pvp}'s discipline --- get the estimand right before admiring the estimator --- has teeth
here.

\textbf{Third, the deliverable is a policy, not a number.} The business buys a \emph{targeting rule}:
a function that says, for each customer, mail or do not mail. Its value is measured in euros, on the
customers it contacts, and \cref{sec:up:euro} measures it.

% =====================================================================================
\section{The data-generating model}
\label{sec:up:dgm}

No causal method can be graded on real data, because on real data the individual effect
$\tau_i$ is never observed for anyone --- that \emph{is} the problem. So the method is first
graded against a world whose answer is known by construction. \Cref{sec:up:valid} is where that
bargain is cashed, and it is the only reason this chapter can say the word ``recovery'' at all.

Each of the \nbOneNCustomers{} customers carries \nbOneNFeatures{} features: recency
$R \sim U(5, 330)$ days since the last purchase, frequency $F \sim \text{Poisson}(5)$, monetary
value $M \sim U(15, 180)$, tenure $\sim U(1, 60)$, and an engagement score
$E \sim \text{Beta}(2, 3)$. With $\sigma(z) = 1/(1 + e^{-z})$ the logistic function, the
simulator plants a baseline spend, a heterogeneous effect, and --- crucially --- a rule that
decides who was emailed in the first place:
%
\begin{align}
  \mu_0(x) &= \max\big\{0,\ 20 + 1.4\,F + 0.25\,M + 12\,E - 0.03\,R\big\},
    \label{eq:up:mu}\\[2pt]
  \tau(x) &= \underbrace{42\, e^{-\left(\frac{R-120}{70}\right)^{2}}\,
                          \sigma(F - 3)\, \sigma\big(9(E - 0.28)\big)}_{\text{persuadables:
                          engaged, mid-recency}}
            \;-\; \underbrace{24\, \sigma\big(9(0.22 - E)\big)\,
                          \sigma\big(1.2(F - 6)\big)}_{\text{sleeping dogs:
                          frequent but disengaged}},
    \label{eq:up:tau}\\[2pt]
  e(x) &= \sigma\Big(1.8\,(E - 0.4) \;+\; 1.2\,\tfrac{150 - R}{150} \;+\; 0.15\,(F - 5)\Big),
          \qquad T \sim \text{Bernoulli}\big(e(x)\big),
    \label{eq:up:prop}\\[2pt]
  Y &= \mu_0(x) \;+\; \tau(x)\, T \;+\; \varepsilon,
       \qquad \varepsilon \sim \mathcal N\!\left(0,\ 8^{2}\right).
    \label{eq:up:y}
\end{align}
%
The constants above are \emph{definitions} of the simulated world, not estimated results ---
\crefrange{eq:up:mu}{eq:up:y} are their complete specification, and \cref{tab:up:params} collects them
for reference. Every number the rest of the chapter reports is measured against them.
\input{build/tables/nb01_params}
Four things are planted here on purpose, and each one is load-bearing.

\textbf{The effect is heterogeneous, and it changes sign.} \Cref{eq:up:tau} is a positive
Gaussian ridge in recency times two logistic gates in frequency and engagement --- the
persuadables --- minus a second surface that switches on for customers who buy often but engage
little: the sleeping dogs. Across this list, \pc{\nbOneSharePersuadable} of customers have
$\tau_i > 0$ and \pc{\nbOneShareSleeping} have $\tau_i < 0$. (The effect is a continuous function,
so nobody is \emph{exactly} zero: the sure things and the lost causes of \cref{sec:up:decision} are
the mass piled near it, not a third share. The euro-relevant subset is smaller than either ---
only \pc{\nbOneShareProfitable} of the list clears the \eur{\nbOneCost} cost of being reached.)
The planted effects run from \eur{\nbOneTauLo} to \eur{\nbOneTauHi}, and their average --- the
number every estimator in \cref{sec:up:classical} is graded against --- is \eur{\nbOneTrueAte}
(\cref{fig:up:types}, left).

\textbf{The confounding is visible in the algebra.} Engagement and recency appear in
\cref{eq:up:prop}, which sets the probability of having been emailed, \emph{and} in
\cref{eq:up:mu}, which sets baseline spend. Marketers historically emailed the engaged and the
recently active --- and those customers spend more anyway. That double role is the whole reason
the naive comparison fails, and \cref{fig:up:types} (right) shows it: the probability of having
been emailed climbs steeply with engagement. In the \emph{randomized} twin used as a control
throughout, \cref{eq:up:prop} is replaced by $e(x) = \tfrac12$ and the confounding vanishes,
with everything else --- including the planted $\tau(x)$ --- held identical.

\textbf{Tenure is a red herring.} It is drawn, it is handed to every estimator, and it enters
neither \cref{eq:up:mu} nor \cref{eq:up:tau}. A good learner must discover that it does not
matter.

\textbf{The noise is large relative to the signal.} $\varepsilon$ has standard deviation
\eur{8} against effects that mostly sit in single digits. Uplift is a
low-signal-to-noise problem by nature --- one is subtracting two noisy surfaces to recover a
small difference --- and that fact, not any philosophical preference, is what decides which
estimator wins in \cref{sec:up:compare}.

\paragraph{A worked handful.} The four types of \cref{sec:up:decision} stop being abstract the
moment \cref{eq:up:tau} is evaluated at a few concrete customers. Take five spanning them --- a
persuadable, a sure thing, a lost cause, a sleeping dog, and a \emph{straddler} whose effect sits
right on the cost line. For customer $i$ the ideal decision is a single comparison,
%
\begin{equation}
  \pi^\star(x_i) \;=\; \mathbf 1\big[\,\tau_i - c \;>\; 0\,\big],
  \qquad c = \eur{\nbOneCost},
  \label{eq:up:oracle}
\end{equation}
%
mail exactly when the planted effect clears the cost of the send, and the euros that send earns are
$\tau_i - c$ itself. Reading $\tau(x)$ off \cref{eq:up:tau} makes the split visible
(\cref{tab:up:types}): the persuadable clears the line and is mailed; the lost cause and the sure
thing sit near zero and are held (the sure thing would have bought anyway, so the discount is margin
given away); and the sleeping dog has $\tau_i - c < 0$, so mailing it does not merely waste $c$ ---
it \emph{subtracts} $c - \tau_i$ from the total. The one type no point comparison can settle is the
straddler, whose effect sits close enough to $c$ that its sign is genuinely uncertain: the customer
\cref{sec:up:euro} pays to test rather than to guess.

\input{build/tables/nb01_worked_types}

\input{build/floats/nb01_types}

\begin{remark}[The confounder is the interesting one, not the noise]
$\hat e(x)$ and $\mu_0(x)$ move together here: customers were emailed \emph{because} they looked
likely to spend. This pattern has a name --- \emph{targeted selection} \citep{hahn2020} --- and it
is the precise regime BCF was designed for. \Cref{sec:up:compare} tests whether that design earns
anything when the confounders are, as here, fully observed.
\end{remark}

% =====================================================================================
\section{Identification}
\label{sec:up:ident}

\subsection*{The estimand, in three rungs}

In the potential-outcome notation of \citet{rubin1974} and \citet{imbens2015}, write $Y_i(1)$ for
customer $i$'s spend if emailed and $Y_i(0)$ for their spend if not. Only one is ever observed ---
the \emph{fundamental problem of causal inference} \citep{holland1986} --- and the individual
effect is their difference:
%
\begin{equation}
  \tau_i \;=\; Y_i(1) - Y_i(0).
  \label{eq:up:ite}
\end{equation}
%
Three estimands are built on \cref{eq:up:ite}, and this chapter needs all three:
%
\begin{align}
  \text{ATE} &\;=\; \E\big[\,Y(1) - Y(0)\,\big]
    &&\text{``does emailing lift spend overall?''} \label{eq:up:ate}\\
  \tau(x) &\;=\; \E\big[\,Y(1) - Y(0) \;\big|\; X = x\,\big]
    &&\text{``the effect for customers \emph{like} $x$'' --- the CATE} \label{eq:up:cate}\\
  V(\pi) &\;=\; \E\big[\,(\tau(X) - c)\,\pi(X)\,\big]
    &&\text{``the euros a targeting rule $\pi$ makes''} \label{eq:up:policy}
\end{align}
%
where $\pi(x) \in \{0, 1\}$ is the policy --- mail or do not mail --- and $c = \eur{\nbOneCost}$
is the cost of a contact. \Cref{eq:up:ate} is what \cref{sec:up:decision} rejected as the
decision object. \Cref{eq:up:cate} is what the rest of the chapter estimates.
\Cref{eq:up:policy} is what the business is actually buying, and it is a \emph{functional} of
$\tau(\cdot)$, not a parameter of any model --- which is why \Cref{ch:pvp}'s remark about
propagating uncertainty through nonlinear functionals comes due here, in euros.

Note what \cref{eq:up:policy} rewards. It does not reward getting $\tau(x)$ right; it rewards
getting the \emph{sign of $\tau(x) - c$} right, customer by customer, and above all it rewards
\emph{ranking} customers correctly, so that the ones contacted first are the ones worth
contacting. A model can be badly wrong about magnitudes and still make almost all the money, and
one in this chapter is exactly that (\cref{fig:up:flatten}). Keep the distinction; it decides
which diagnostic matters.

\subsection*{The assumptions}

Recovering \cref{eq:up:cate} from data in which treatment was \emph{not} randomized rests on
three assumptions. Each is stated, and each is either tested later in the chapter or explicitly
declared untestable. An assumption with no accompanying test is a wish.

\begin{assumption}[Unconfoundedness]\label{as:up:unconf}
Given the measured features, who was emailed is as good as random:
%
\begin{equation}
  \big\{\,Y(0),\, Y(1)\,\big\} \;\perp\!\!\!\perp\; T \;\;\big|\;\; X .
  \label{eq:up:unconf}
\end{equation}
%
This is the big one, and it is \textbf{untestable from the data} \citep{rosenbaum1983}. It holds
by construction in \cref{eq:up:prop}, because the simulator built the assignment rule out of $X$
and nothing else; it holds by \emph{design} in a randomized experiment; and it is an act of faith
everywhere else. It is not tested in \cref{sec:up:valid} --- it cannot be. Instead
\cref{sec:up:wrong} asks the only question available: how strong would a confounder have to be,
had one been missed, to overturn the decision?
\end{assumption}

\begin{assumption}[Overlap / positivity]\label{as:up:overlap}
Every kind of customer had some chance of being emailed and some chance of not:
$0 < e(x) < 1$ for all $x$ in the support, where $e(x) = \Prob(T = 1 \mid X = x)$ is the
\emph{propensity score}. Where a whole region of feature space was always emailed there is
nobody to compare against, and any estimator is silently extrapolating rather than comparing.
This one \emph{is} testable, and it is tested first, before anything is fitted
(\cref{fig:up:overlap}): only \nbOneWeakN{} of \nbOneNCustomers{} customers
(\pc{\nbOneWeakPct}) sit in the $<\!0.05$ or $>\!0.95$ tails.
\end{assumption}

\begin{assumption}[SUTVA]\label{as:up:sutva}
One customer's email does not change another's spend --- no word-of-mouth, no shared household,
no forwarding of the discount code --- and there is only one version of ``the email''
\citep{rubin1980, imbens2015}. In a mailing list this is usually benign and occasionally
catastrophic: a discount code that leaks to a coupon site violates it wholesale, and no
diagnostic in this chapter would notice.
\end{assumption}

Under \crefrange{as:up:unconf}{as:up:sutva} the \emph{backdoor adjustment} \citep{pearl2009}
identifies the CATE as a contrast of two observable regression surfaces:
%
\begin{equation}
  \tau(x) \;=\; \underbrace{\E[\,Y \mid T = 1,\ X = x\,]}_{\textstyle \mu_1(x)}
             \;-\; \underbrace{\E[\,Y \mid T = 0,\ X = x\,]}_{\textstyle \mu_0(x)} .
  \label{eq:up:backdoor}
\end{equation}
%
\Cref{eq:up:backdoor} is the whole causal content of this chapter, and it is worth saying plainly
what it is \emph{not}: it is not a model, not a prior, and not a sampler. Everything after it is
statistics. The confounding is defeated by \emph{comparing emailed against not-emailed customers
who look alike}, and every estimator below --- the classical ones and the Bayesian one --- is a
different way of estimating the same two conditional means. This is \Cref{ch:pvp}'s central
discipline, restated: \textbf{the graph decides what enters the model; the machinery only prices
the coefficients}.

% =====================================================================================
\section{The classical baseline}
\label{sec:up:classical}

Before any sampler, what does a competent analyst produce \emph{without} one? The answer is: very
nearly everything. This section builds a complete uplift analysis --- the average effect with
honest standard errors, a per-customer effect, a ranking, and even a per-customer interval ---
with no likelihood, no prior and no chain anywhere in it. That is the bar \cref{sec:up:bayes}
must clear, and it is deliberately not a strawman.

\subsection*{The average effect, six ways --- and the first is a trap}

The simplest estimator, the difference in means
$\bar y_{\text{emailed}} - \bar y_{\text{not emailed}}$ with a Welch standard error, is
\emph{exactly right} under randomization. Here assignment came from the selection rule
\cref{eq:up:prop}, so it is exactly wrong, and confidently so. It returns \eur{\nbOneNaiveEst}
against a planted \eur{\nbOneTrueAte} --- an error of \nbOneNaiveErr{} euros, nearly double the truth
--- inside a 90\,\% interval of $[\eur{\nbOneNaiveLo},\ \eur{\nbOneNaiveHi}]$ that misses it. This is
the shape of \emph{precisely wrong}.

Run the identical estimator on a \textbf{randomized twin} of the world --- same customers, same
planted $\tau(x)$, \cref{eq:up:prop} replaced by a coin flip --- and it returns \eur{\nbOneRctEst}.
Same estimator, same arithmetic, opposite verdicts: \textbf{what broke was the identification, not the
estimator}.

Adjustment repairs it, three ways (\cref{tab:up:ate}, whose sixth row carries the Bayesian posterior
of \cref{sec:up:bayes} so the average reads in one place):

\begin{itemize}[leftmargin=*]
  \item \textbf{OLS with the features}, $y \sim T + x$, with an HC1 heteroskedasticity-robust
        covariance \citep{white1980, mackinnon1985}. Spend variance differs by arm and by customer
        type, so an iid standard error would be a fiction. It gives \eur{\nbOneOlsEst}.
  \item \textbf{OLS, Lin-interacted}, $y \sim T \times (x - \bar x)$. Centring the covariates makes
        the coefficient on $T$ the ATE; giving each arm its own slopes keeps the estimator consistent
        for the ATE even when $\tau(x)$ is heterogeneous --- \textbf{under randomized assignment}, and
        whatever the functional form \citep{lin2013}. That scope matters here. On this
        \emph{observational} panel the adjustment carries the identification, and \cref{eq:up:mu} is
        not linear (it clips at zero), so a linear adjustment only approximates and nothing promises
        consistency. The residual error is not assumed away; the next subsection measures it, at
        \nbOneLinPanelBias{} euros and \nbOneLinPanelSes{} standard errors from zero. Plain additive
        OLS is worse: it assumes one common effect and, under real heterogeneity, targets a
        variance-weighted average --- a bias the interaction removes for \nbOneNFeatures{} extra
        coefficients. It gives \eur{\nbOneLinEst}.
  \item \textbf{AIPW}, the doubly-robust estimator, at \eur{\nbOneAipwEst} with a
        \nbOneNBootAipw-resample bootstrap interval of
        $[\eur{\nbOneAipwLo},\ \eur{\nbOneAipwHi}]$.
\end{itemize}

AIPW deserves its algebra, because ``doubly robust'' is otherwise a slogan. It averages one
\emph{influence-function} score per customer --- that customer's contribution to the average effect,
built so an error in either model washes out:
%
\begin{equation}
  \hat\psi_i \;=\; \underbrace{\hat\mu_1(x_i) - \hat\mu_0(x_i)}_{\text{outcome-model guess}}
    \;+\; \underbrace{\frac{T_i\big(y_i - \hat\mu_1(x_i)\big)}{\hat e(x_i)}
        \;-\; \frac{(1 - T_i)\big(y_i - \hat\mu_0(x_i)\big)}{1 - \hat e(x_i)}}_{
        \text{inverse-propensity-weighted residual corrections}},
  \qquad
  \widehat{\text{ATE}} \;=\; \frac{1}{n}\sum_{i=1}^{n} \hat\psi_i .
  \label{eq:up:aipw}
\end{equation}
%
If the outcome models are right, the residuals $y_i - \hat\mu_{T_i}(x_i)$ are mean-zero, the
correction terms vanish on average, and the first term stands. If instead the propensity is right,
the inverse-weighted residuals repair the bias the outcome models leave. \emph{Either} model being
correct suffices; only both wrong breaks it \citep{robins1994}. The estimator here is
\textbf{cross-fitted} \citep{chernozhukov2018}: the boosted-tree outcome models and the logistic
propensity are fitted on $K-1$ folds and scored out of fold, so no learner grades its own homework
and \cref{eq:up:aipw}'s residuals are not shrunk by own-observation bias.

\input{build/tables/nb01_ate}

\subsection*{Bias or luck? A distinction the table cannot make}

One panel cannot separate a systematically biased estimator from an unlucky draw. So the same
three estimators are re-run on \nbOneNPanels{} fresh panels of \emph{each} regime, and the mean
error against each panel's own planted ATE is measured.

The naive estimator on confounded panels is off by \nbOneNaivePanelBias{} euros on average ---
\nbOneNaivePanelSes{} standard errors of the mean from zero. It is not unlucky; it is
\textbf{wrong}, by roughly the same amount every time. The identical estimator on
\emph{randomized} panels sits at \nbOneRctPanelBias{} euros, \nbOneRctPanelSes{} standard errors
from zero. Nothing changed but how customers were selected.

The adjusted arm's residual error is \nbOneLinPanelBias{} euros --- \nbOneLinPanelSes{} standard
errors from zero, so it is \nbOneResidVerdict{} --- roughly \nbOneBiasShrink$\times$ smaller than
the naive arm's, and \pc{\nbOneLinResidPct} of the effect being measured. Adjustment
\emph{repairs} the identification; it does not \emph{anoint} the estimate.

\begin{remark}[The yardstick, stated rather than defaulted]
A mean error counts as bias when it is far from zero relative to \emph{its own standard error},
$s/\sqrt{n_{\text{panels}}}$ --- not relative to the panel-to-panel standard deviation $s$, which
measures something else entirely (how much one draw jumps around) and is a bar lax enough to wave
a genuine four-sigma bias through as noise. The notebook applies the first, and prints how many
standard errors from zero each arm sits.
\end{remark}

\subsection*{How a CATE model is graded}

Four numbers recur from here to the end of the chapter, and all four are defined once, now.

\begin{description}[leftmargin=*, style=nextline]
  \item[PEHE --- Precision in Estimation of Heterogeneous Effects.] The root-mean-square error of
        the estimated effect against the true one,
        $\text{PEHE} = \sqrt{\tfrac1n \sum_i (\hat\tau(x_i) - \tau_i)^2}$, in euros. Lower is
        better. It is the CATE's answer to RMSE --- and, being a comparison against $\tau_i$, it
        is available \emph{only in a simulator}. \Cref{sec:up:valid} says what replaces it in
        production.
  \item[Correlation.] The correlation between $\hat\tau(x_i)$ and $\tau_i$ across customers: a
        crude but honest summary of whether the model has the \emph{shape} of the effect surface.
  \item[Qini / uplift curve and AUUC.] Sort customers by \emph{predicted} effect, contact them
        from the top down, and plot the cumulative \emph{true} effect captured against the
        fraction contacted. Random targeting traces the diagonal; an oracle that knows every
        $\tau_i$ traces the upper envelope. \textbf{AUUC} is the area between the model's curve
        and the diagonal, normalised by the oracle's, so $1.0$ is a perfect ranking and $0$ is no
        better than random \citep{radcliffe2007, radcliffe2011, gutierrez2017}. It scores the only
        thing \cref{eq:up:policy} consumes: the \emph{order}.
  \item[Coverage.] The share of customers whose true $\tau_i$ falls inside their nominal 90\,\%
        interval. It is reported overall \emph{and} by decile of the true effect, because a model
        can be well calibrated on average and badly overconfident exactly where the money is ---
        and one in this chapter is.
\end{description}

\subsection*{Heterogeneous effects, with no Bayes anywhere}

The \textbf{T-learner} \citep{kunzel2019} is \cref{eq:up:backdoor} done twice and subtracted:
%
\begin{equation}
  \hat\tau(x) \;=\; \hat\mu_1(x) - \hat\mu_0(x),
  \qquad \hat\mu_t = \text{a regression of $y$ on $x$, fitted on the arm } T = t .
  \label{eq:up:tlearner}
\end{equation}
%
Two scikit-learn fits and a subtraction. The base learner is a gradient-boosting machine
\citep{friedman2001} --- 200 trees, depth 3, one fixed and unremarkable configuration used
identically everywhere in the chapter. The \textbf{S-learner} is its rival recipe: fit
\emph{one} model on $(x, T)$ together and read the effect as
$\hat f(x, 1) - \hat f(x, 0)$. And a third fit, the T-learner recipe with a \emph{linear} model
per arm, is the control experiment for functional form.

\Cref{tab:up:cate} grades all three against the planted $\tau(x)$. The GBM T-learner reaches an
AUUC of \nbOneGbmTAuuc{} of the oracle's ranking, from two scikit-learn calls and a subtraction.
\textbf{Anyone who says a posterior is required to do uplift is selling something.} What is
required is (i) the identification argument of \cref{sec:up:ident} and (ii) a base learner
flexible enough for the true $\tau(x)$ --- and the linear row shows the second is not free: the
same recipe with a straight line per arm collapses to AUUC \nbOneOlsTAuuc, because
\cref{eq:up:tau} is a smooth non-linear ridge and a plane cannot bend to it.

\input{build/tables/nb01_cate}

The S-learner row is the one to read carefully, because its failure is famous and usually
mis-stated. Its single shared surface lets the regulariser shrink the treatment column's
influence toward zero, so it \textbf{flattens} $\tau(x)$: its PEHE is \eur{\nbOneGbmSPehe}
against the T-learner's \eur{\nbOneGbmTPehe}. And yet with \emph{this} base learner it costs
essentially nothing on the ranking (AUUC \nbOneGbmSAuuc{} against \nbOneGbmTAuuc). A uniformly
squashed $\tau(x)$ can still put customers in the right \emph{order}. \Cref{sec:up:wrong} repeats
the experiment with a BART base learner and finds the damage lands somewhere else entirely ---
which is the real lesson: \emph{which} axis the S-learner breaks is a property of the base
learner, not of the letter S.

\subsection*{Where the classical arm runs out of road --- measured, not asserted}

\Cref{eq:up:tlearner} returns a \emph{point} $\hat\tau(x)$ and nothing else. No per-customer
interval falls out of it. The classical escape hatch is the nonparametric bootstrap
\citep{efron1979}: resample the customers, refit both arms, repeat \nbOneNBootCate{} times, and
read off a band per customer. It is cheap, so the notebook computes it and then grades it against
the planted truth. Its 90\,\% band covers \pc{\nbOneBootCov} of customers overall, at a mean width
of \eur{\nbOneBootWidth} --- and \pc{\nbOneBootTopDec} on the top decile of the true effect
against \pc{\nbOneBootBotDec} on the bottom. Broadly honest in the middle, fraying at both
extremes, which is precisely where a targeting rule lives.

But grade the \emph{object}, not only the number. A bootstrap band describes the sampling
variability of the \emph{estimator} --- where $\hat\tau(x)$ would land if the customer list were
redrawn. It is centred on $\hat\tau(x)$, not on $\tau_i$. Every replicate refits the same learner
on the same information, so a systematically shrunken estimate stays shrunken in all
\nbOneNBootCate{} of them, and the band tightens \textbf{around the wrong value} without a hint of
complaint. \Cref{sec:up:euro}'s rule needs $\Prob(\tau_i > c)$ --- a probability about
\emph{this customer's effect} --- and \Cref{ch:pvp} settles that the classical arm has no
such object to offer. Whether that distinction costs real money, or is a philosopher's quibble,
this chapter will not assert. \Cref{sec:up:compare} measures it, in euros.

% =====================================================================================
\section{The Bayesian model}
\label{sec:up:bayes}

The Bayesian arm is a \textbf{Bayesian Causal Forest} \citep{hahn2020}, which is BART
\citep{chipman2010, hill2011} arranged for a causal contrast. BART is a sum-of-trees model with
a regularizing prior that keeps each tree shallow, so that no single tree explains much and the
ensemble is an average of many weak, mutually-correcting learners; the object it returns is a
\emph{posterior} over functions, not a fitted function. BCF makes two changes to it, and both are
claims about the world rather than about computation:
%
\begin{align}
  y_i \;\big|\; \mu_i, \sigma
    &\;\sim\; \mathcal N\big(\mu_i,\ \sigma^{2}\big),
    \qquad
    \mu_i \;=\; f_{\text{prog}}\big(x_i,\ \hat e(x_i)\big) \;+\; f_{\tau}(x_i)\, T_i ,
    \label{eq:up:bcf}\\
  f_{\text{prog}} &\;\sim\; \text{BART}(m \text{ trees}),
  \qquad
  f_{\tau} \;\sim\; \text{BART}(m/2 \text{ trees}),
  \qquad
  \sigma \;\sim\; \text{HalfNormal}(15).
  \label{eq:up:bcfprior}
\end{align}
%
The \emph{prognostic} surface $f_{\text{prog}}$ models baseline spend; the \emph{treatment}
surface $f_{\tau}$ is the CATE, and its posterior \emph{is} $\hat\tau(x)$ ---
$\hat\tau(x) = f_\tau(x)$, a full posterior distribution per customer, of shape
(draws $\times$ customers).

\textbf{Change one: the propensity goes in as a covariate.} $\hat e(x)$, estimated by logistic
regression, is handed to the prognostic surface. Under targeted selection --- customers emailed
\emph{because} they looked likely to spend --- the propensity and the baseline move together, and
giving the model $\hat e(x)$ lets it soak up exactly that shared structure instead of leaking it
into the treatment surface \citep{hahn2020}. That is the trick BCF is famous for.

\textbf{Change two: the treatment surface gets half as many trees.} This is the
\emph{shrink-toward-homogeneity} prior in its concrete form: heterogeneity is harder to express
than baseline structure, so the model reports it only when the data insist. It is a claim that
effects are more nearly constant than baselines are --- defensible, common, and, on this data,
tested in \cref{sec:up:compare} and found wanting.

One specification note, because the chapter's central negative finding leans on it. \code{pymc-bart}
centres each forest's prior at the mean of the response it is handed, and here the treatment forest is
handed the \emph{outcome}. Its prior therefore sits at the outcome level, not at zero as in the clean
Hahn--Murray--Carvalho form where the $\tau$ prior shrinks toward \emph{no effect}. In
\cref{eq:up:bcfprior} the shrinkage comes from the halved tree count, not from a zero-centred
treatment prior. Every verdict below is scoped to \emph{this} parameterisation of BCF.

Inference is by the particle-Gibbs sampler PGBART \citep{quiroga2022}, implemented in
\code{pymc-bart} on top of PyMC \citep{salvatier2016}: \nbOneBcfDraws{}
post-warmup draws across \nbOneBcfChains{} chains, with $m = \nbOneBcfTrees$ trees on the
prognostic surface. The diagnostics scoped to the model's continuous scalars give
$\hat R = \nbOneBcfRhat$ and a minimum bulk effective sample size of \nbOneBcfEss{}
\citep{vehtari2021}. For a tree sampler these scalars are a weaker signal than they look, so the
convergence verdict is computed and printed rather than eyeballed: on this run it is
\textbf{\nbOneBcfRhatVerdict}. The band that produces that word --- and why $\hat R$ is a poor proxy
for what the chapter actually reads off the model, the treatment surface --- is set out with a
worked example in \cref{app:up:tol}.

Two facts about \code{pymc-bart} make that caution load-bearing. First, the real backstop for
convergence is not $\hat R$ but the \nbOneNSeeds-seed refit of \cref{sec:up:valid}, which grades the
\emph{estimates} rather than the sampler and finds them \nbOneStabVerdict. Second, and more
uncomfortable, the sampler is not seed-stable in this environment: two full runs of identical source
have moved this model's PEHE by more than a fifth. So no BART-derived number is ever typed into this
chapter --- every one is injected from the run that produced the PDF in front of you --- and every
comparison between two of them is judged against a \emph{measured} tolerance rather than by eye. That
tolerance is constructed in \cref{app:up:tol}, and it is why \cref{sec:up:compare} is built the way
it is.

\subsection*{Model criticism, before any effect is read off}

The posterior-mean ATE is \eur{\nbOneBcfAte}, which lies \nbOneBcfAgrees{} the AIPW 90\,\%
interval of $[\eur{\nbOneAipwLo},\ \eur{\nbOneAipwHi}]$ --- two very different estimators, one
classical and doubly robust, one a Bayesian tree ensemble, agreeing on the average. That is the first credibility check, and it is banked before anything harder is
attempted.

The second is a genuine posterior predictive check (\cref{fig:up:ppc}): replicate datasets are
drawn from the \emph{fitted} model --- each from its own draw's mean surface and noise scale ---
and their spend distribution laid over the real one. \pc{\nbOnePpcCov} of observed spend falls
inside the replicates' 5--95\,\% range, against the 90\,\% a calibrated model would give:
\nbOnePpcVerdict. Split by arm --- the sharper version, because $\tau(x)$ \emph{is} the difference
between the two arms' surfaces --- the same check gives \pc{\nbOnePpcCovTreated} among the emailed
and \pc{\nbOnePpcCovControl} among the not-emailed, so the pooled pass is not concealing a one-arm
misfit. A model that cannot generate the data it was trained on has no business being read for
effects, however tight its intervals.

Bank that pass with its scope attached, because it is a weak test and the chapter does not pretend
otherwise. A flexible sum-of-trees mean surface with a free noise scale reproduces the marginal of
spend almost by construction; a check the model can barely fail is not evidence that it has the
\emph{heterogeneity} right, and no marginal check --- pooled or arm-split --- can see per-customer
error in $\tau(x)$. The strict test is the truth-graded coverage by decile of
\cref{sec:up:valid}, which the simulator makes possible and which this model does \emph{not} pass.

\input{build/floats/nb01_overlap}
\input{build/floats/nb01_ppc}

% =====================================================================================
\section{Diagnostics and validation}
\label{sec:up:valid}

\subsection*{The problem: you cannot score a CATE model directly}

Here is the difficulty that makes uplift modelling strange, and it is worth stating before the
diagnostics rather than after. For a spend forecast, validation is arithmetic: hold out customers,
compare prediction with realisation. For a \emph{causal effect}, there is nothing to compare
against. $\tau_i$ is not a missing label that better data collection would supply; it is a
difference between one observed future and one that never happened. \textbf{No test set contains
it, ever.}

Two responses, and this chapter uses both.

\textbf{The simulator's response} is to plant $\tau_i$ and grade against it. That is what PEHE,
the correlation, the coverage and the recovery scatter of \cref{fig:up:recovery} all do, and it
is the only reason this chapter can make a statement like ``the intervals under-cover the largest
effects'' at all. \textbf{Say plainly what this costs: PEHE is a luxury of the simulator.} On the
real customer list it does not exist, cannot be approximated, and any vendor quoting one is
quoting a number from a world they built themselves.

\textbf{The production response} is the Qini curve, and its logic is worth spelling out because
it looks like magic. Rank customers by predicted uplift. Take the top $k$. Within that top-$k$
slice --- \emph{if treatment was randomized} --- the emailed and the not-emailed are exchangeable,
so their difference in mean outcome is an unbiased estimate of the average effect \emph{on that
slice}. Sweep $k$ and one traces out the cumulative incremental gain from targeting the top $k$,
against what random targeting would have delivered. No counterfactual is needed anywhere:
the estimator borrows the randomization to construct, at each $k$, the comparison the individual
customer cannot supply. That is why AUUC survives contact with production and PEHE does not.

And it is worth being equally plain about what Qini \emph{cannot} see. It grades the
\emph{ranking}, not the magnitudes, so a model that flattens every effect by a factor of five
scores identically to the correct one --- \cref{fig:up:flatten} is that fact, drawn. It says
nothing about whether $\hat\tau(x)$ clears the cost \emph{in euros}, which is what
\cref{eq:up:policy} actually needs. It is noisy in the top decile, where the slices are smallest
and the money is. And it requires randomized data: on an observational panel the top-$k$
difference in means is confounded exactly as \cref{tab:up:ate}'s first row is, and the curve
inherits that bias wholesale.

\subsection*{Recovery, calibration and ranking}

\Cref{fig:up:recovery} and \cref{fig:up:qini} are the payoff of having a known truth. The
posterior-mean $\hat\tau(x)$ correlates with the planted $\tau(x)$ at \nbOneBcfCorr, with a PEHE
of \eur{\nbOneBcfPehe} against effects that run from \eur{\nbOneTauLo} to \eur{\nbOneTauHi}. It
ranks customers at an AUUC of \nbOneBcfAuuc{} of the oracle's, sorting the list into a clean
descending staircase from \eur{\nbOneDecTopTrue} of true effect in the top predicted decile to
\eur{\nbOneDecBotTrue} in the bottom.

The coverage is where the interesting problem is. Overall, the nominal 90\,\% credible interval
covers \pc{\nbOneBcfCovNinety} of customers --- on its nominal target, and unremarkable. Split by
decile of the \emph{true} effect and the average dissolves into a shape: coverage sags at the
extremes, and on the \textbf{top} decile --- the customers a targeting rule is about to spend
money on --- it falls to \pc{\nbOneBcfTopDec}. Read the reliability panel of \cref{fig:up:recovery}
alongside it and the mechanism is visible, and it is a failure of \emph{location} rather than of
width: the shrink-toward-homogeneity prior of \cref{eq:up:bcfprior} pulls the biggest effects
\emph{down}, and the interval, centred on the shrunken value, no longer reaches back to the truth.
The widths are printed rather than assumed: on that top decile BCF's mean 90\,\% credible interval
is \eur{\nbOneBcfTopdecWidth}, \nbOneTopdecWidthVerdict{} the classical bootstrap band's
\eur{\nbOneBootTopdecWidth}, so what carries the miss is the shrunken \emph{centre} and not an
appeal to sharpness. \textbf{The model is most confident, and most wrong, exactly where the euros
are.}

\input{build/floats/nb01_recovery}
\input{build/floats/nb01_qini}

\subsection*{Stability --- and an honest word about what ``stable'' means here}

One panel is one draw of the world. Refitting the whole pipeline on \nbOneNSeeds{} fresh
simulations gives an ATE recovery bias of \nbOneStabBias{} euros with a seed-to-seed standard
deviation of \eur{\nbOneStabSd}: \nbOneStabVerdict.

But those \nbOneNSeeds{} refits also measure something the chapter cannot politely ignore. Across
them, PEHE ranges from \eur{\nbOneStabPeheLo} to \eur{\nbOneStabPeheHi} and AUUC from
\nbOneStabAuucLo{} to \nbOneStabAuucHi. Some of that is fresh data and some of it is the sampler,
and \cref{sec:up:bayes} is the reason the two cannot be cleanly separated here. The operational
consequence is a rule this chapter obeys without exception: \textbf{no BART-derived number is ever
compared with another by eye}. Every head-to-head below is judged against a measured tolerance ---
twice the paired-bootstrap standard error of the gap, plus the sampler's own fit-to-fit wobble
wherever a BART refit sits on one side of the comparison --- so a gap smaller than its own noise
prints as ``too close to call'' rather than as a winner. \Cref{app:up:tol} constructs that
tolerance and works an example on the PEHE head-to-head.

% =====================================================================================
\section{Which model do we trust for the targeting rule?}
\label{sec:up:compare}

\input{build/tables/nb01_compare}

The rule the business ships thresholds a per-customer probability, $\Prob(\tau_i > c \mid D)$, so
this section asks one question and answers it forward: \emph{which model should the decision trust
to produce that number and to rank customers?} The two arms are put on the same estimands, the
places they agree are banked, the places they differ are audited ingredient by ingredient, and the
choice of decision model falls out at the end.

\subsection*{First, the average: the two arms agree}

\Cref{tab:up:ate} and \cref{tab:up:compare} put the two arms on the same estimands. On the ATE,
the classical estimators and the BCF posterior land on the same effect --- the posterior mean is the
last row of \cref{tab:up:ate}, printed there with its credible interval so that the agreement can be
read off one table rather than assembled from two --- all of them within a euro
of the planted \eur{\nbOneTrueAte}. \textbf{The Bayesian layer did not find a better number, and
it could not have.} With \nbOneNCustomers{} customers, fully observed confounders and weak priors,
a correctly adjusted regression is already consistent and there was nothing left on the table.
This is Bernstein--von Mises --- \Cref{ch:pvp}'s proposition --- arriving on schedule, and it is the
ordinary case in this book. If you came expecting ``going Bayesian'' to fix a biased estimate,
\cref{sec:up:classical} is the correction: what fixed the bias was adjusting for the
\nbOneNFeatures{} features, and OLS does that for free.

\subsection*{On the per-customer effect they separate --- so we audit the ranking}

On the per-customer effect the two arms do come apart, and \cref{tab:up:compare} records it. The
head-to-head, judged against the measured tolerance, names \textbf{\nbOneHhPeheWin} on PEHE
(\nbOneHhPeheA{} against \nbOneHhPeheB, a gap of \nbOneHhPeheGap{} against a tolerance of
\nbOneHhPeheTol) and \textbf{\nbOneHhAuucWin} on AUUC (\nbOneHhAuucA{} against \nbOneHhAuucB). A
winner, though, is not yet a reason to trust one model over the other. Before the decision leans on
BCF's ranking we have to know \emph{what} in BCF produced it.

BCF differs from the classical GBM T-learner in \emph{four} ways at once: (i) it averages over a
\emph{posterior} of tree ensembles instead of committing to a single point fit; (ii) it applies a
shrink-toward-homogeneity prior to the treatment surface --- in this implementation, half as many
trees on it; (iii) it feeds in the propensity $\hat e(x)$; and (iv) it fits \emph{one} joint
additive surface, $f_{\text{prog}} + f_\tau T$, where the T-learner fits two independent forests. A
two-horse race cannot say which of the four did the work, and the flattering guess --- that the
credit belongs to (ii) and (iii), the specifically \emph{causal} machinery a vendor sells --- is
exactly the one a buyer should be most sceptical of. So we isolate each ingredient before trusting
the ranking.

\subsection*{Due diligence: the attribution ladder}

Fit a third model: a \textbf{BART T-learner}. Two arms, subtract --- \cref{eq:up:tlearner} exactly ---
but with a sum-of-trees \emph{posterior} as the base learner in place of the gradient-boosting point
fit. It has ingredient (i) and none of (ii), (iii), (iv). That gives a ladder. Read each rung by what
it moves:

\begin{center}\small
\begin{tabular}{@{}lll@{}}
\toprule
\textbf{Rung} & \textbf{What changes} & \textbf{What it prices} \\
\midrule
(1) GBM T-learner $\to$ (2) BART T-learner
  & point fit $\to$ posterior
  & \textbf{the base learner} \\
(1) GBM T-learner $\to$ (1b) bagged GBM
  & one fit $\to$ the mean of many
  & \textbf{averaging alone} \\
(2) BART T-learner $\to$ (3) BCF
  & $+$ shrink prior, $+$ propensity,
  & \textbf{BCF's package}, \\
  & $+$ one joint surface
  & not any one ingredient \\
\bottomrule
\end{tabular}
\end{center}

Rung (2) $\to$ (3) is a \emph{bundle}. Three ingredients move together there, and no verdict below can
name which one is responsible; the ladder credits or debits BCF \emph{as specified}. Isolating the
prior alone would take a fourth fit --- BCF with an unhalved treatment forest, the propensity still
fed in --- and this chapter does not run it. What points at the prior \emph{specifically} is
\cref{fig:up:recovery}'s reliability panel, where the largest effects are visibly flattened.

\input{build/tables/nb01_ladder}

\Cref{tab:up:ladder} is the result. Every gap is judged against the measured tolerance of
\cref{app:up:tol}: twice its paired-bootstrap standard error, plus, on the BCF rung, the sampler's own
fit-to-fit wobble that the paired bootstrap cannot see. That wobble is real. The BCF in
\cref{tab:up:ladder} is an equal-compute \emph{refit} of the model \cref{tab:up:compare} reports ---
same specification, same customers, a different seed --- and the two fits disagree by
\eur{\nbOneBcfSeedPeheWobble} of PEHE and \nbOneBcfSeedCovWobble{} points of overall coverage, purely
from the sampler. That disagreement also reconciles the two tables: they show two fits of one model.
Folding it into the tolerance lets a gap smaller than the sampler's noise print as ``too close to
call''. The verdicts, computed by the notebook and cited rather than re-argued:

\begin{itemize}[leftmargin=*]
  \item \textbf{Rung (1) $\to$ (2), the base learner.} On PEHE the winner is
        \textbf{\nbOneRungBlPeheWin} (\nbOneRungBlPeheA{} against \nbOneRungBlPeheB, a gap of
        \nbOneRungBlPeheGap{} against a tolerance of \nbOneRungBlPeheTol). On AUUC it is
        \textbf{\nbOneRungBlAuucWin} (\nbOneRungBlAuucA{} against \nbOneRungBlAuucB, gap
        \nbOneRungBlAuucGap, tolerance \nbOneRungBlAuucTol).
  \item \textbf{Rung (1) $\to$ (1b), averaging alone.} The same gradient-boosting T-learner,
        averaged over the \nbOneNBootCate{} bootstrap refits the classical arm already computed:
        model averaging with the base learner held fixed and no Bayes anywhere. On PEHE the winner
        is \textbf{\nbOneRungBagPeheWin} (\nbOneRungBagPeheA{} against \nbOneRungBagPeheB); on AUUC,
        \textbf{\nbOneRungBagAuucWin} (\nbOneRungBagAuucA{} against \nbOneRungBagAuucB). This is the
        control that keeps rung (1) $\to$ (2) honest: whatever averaging alone does not deliver, the
        rest of that rung --- a different ensemble, with a regularizing tree prior --- is what did.
  \item \textbf{Rung (2) $\to$ (3), BCF's package.} On PEHE: \textbf{\nbOneRungBcfPeheWin}
        (\nbOneRungBcfPeheA{} against \nbOneRungBcfPeheB, gap \nbOneRungBcfPeheGap, tolerance
        \nbOneRungBcfPeheTol). On AUUC: \textbf{\nbOneRungBcfAuucWin} (\nbOneRungBcfAuucA{} against
        \nbOneRungBcfAuucB). On overall coverage: \textbf{\nbOneRungBcfCovWin}. And on the axis
        that matters most --- coverage of the \emph{top decile of true effect} ---
        \textbf{\nbOneRungBcfTopdecWin} (\nbOneRungBcfTopdecA{} against \nbOneRungBcfTopdecB, gap
        \nbOneRungBcfTopdecGap{} against a tolerance of \nbOneRungBcfTopdecTol).
\end{itemize}

The notebook's summary, from those verdicts: the base learner buys \textbf{\nbOneAttrBaselearner};
BCF's package buys \textbf{\nbOneAttrBcfBuys} and costs \textbf{\nbOneAttrBcfCosts}. In one sentence,
the notebook's: the Bayesian edge over the classical arm is \nbOneAttrVerdict.

State the verdict plainly. What moves the needle is the ingredient nobody sells: \textbf{a
sum-of-trees posterior in place of a single point fit}. Rung (1) $\to$ (1b) keeps that from
over-reaching --- it prices the averaging alone, base learner held fixed --- and what remains of rung
(1) $\to$ (2) is the ensemble itself: many shallow, mutually-correcting trees under a regularizing
prior, integrated rather than fitted. That buys real accuracy when the signal is thin, and it is thin
here: \nbOneNCustomers{} customers, \eur{8} of noise per observation, effects mostly in single digits.
This is the \emph{ordinary} Bayesian gain \citet{chipman2010} and \citet{hill2011} would have
predicted, nothing exotic. BCF's specifically \emph{causal} package adds \textbf{\nbOneAttrBcfBuys} on
top, and \cref{sec:up:dgm}'s remark said why: the propensity trick absorbs targeted selection, and
with the confounders \emph{fully observed} and overlap good, there is nothing left for $\hat e(x)$ to
fix.

The point for the decision is not who \emph{wins}. It is that the ranking's accuracy is real and
\emph{attributable} to the base learner, so the decision can lean on the ranking on its merits and
price BCF's causal package separately. The next two subsections do that pricing: the intervals, then
the euros.

\subsection*{The intervals: BCF is least trustworthy where the rule spends money}

Overall coverage is the least interesting line in \cref{tab:up:compare}: the classical band
under-covers at \pc{\nbOneBootCov}, the posterior lands on \pc{\nbOneBcfCovNinety}, and the
head-to-head returns \textbf{\nbOneHhCovWin}. Read that line as
``the same, to within the noise'' whichever name it prints --- the gap has landed on both sides of
its tolerance across consecutive full runs of identical source, which is a lesson in why this
chapter computes its verdicts instead of typing them.

Split by decile of the true effect and the arms come apart, and this one is not marginal. On the
top decile the winner is \textbf{\nbOneHhTopdecWin} (\pc{\nbOneBootTopDec} against
\pc{\nbOneBcfTopDec}), a gap of \nbOneHhTopdecGap{} against a tolerance of \nbOneHhTopdecTol. The
cause is not a bug, and it is a failure of \emph{location}: shrinkage toward homogeneity pulls the
largest $\tau_i$ down, and the interval --- \eur{\nbOneBcfTopdecWidth} wide on that decile,
\nbOneTopdecWidthVerdict{} the classical band's \eur{\nbOneBootTopdecWidth} --- is centred on the
shrunken value and no longer reaches back to the truth. The classical arm,
having no shrinkage to apply, is less accurate on average and more honest at the extreme.

Now put that next to the ladder, because together they are the sharpest thing in this chapter. The
shrinkage that costs the top decile its coverage belongs to the package the ladder found buying
\nbOneAttrBcfBuys. \textbf{On this world it is close to all cost and no benefit --- and that is
known because it was measured, not because it sounds humble.} A prior is a claim about the world,
and this world declines \emph{this} one --- the halved treatment forest of
\cref{eq:up:bcfprior}, in the parameterisation \cref{sec:up:bayes} disclosed. The ladder debits the
package; what points at the shrinkage in particular is \cref{fig:up:recovery}'s reliability panel,
where the largest effects are visibly flattened.

\subsection*{In euros: does the posterior's rule beat the best classical rule?}

\input{build/tables/nb01_rules}

\Cref{tab:up:rules} is the only grade that pays anyone, and its first two rows are what the
\emph{average} can instruct: email everyone, which destroys \eur{\nbOneTreatAllLoss} and is what the
naive \eur{\nbOneNaiveEst} recommends, or email nobody, which is what the correct \eur{\nbOneTrueAte}
recommends and which earns zero against an oracle prize of \eur{\nbOneOracleProfit}. Every
model-driven rule recovers the great majority of that prize. And then look at the rules that
use their uncertainty: the posterior's honest rule, $\Prob(\tau_i > c) > 0.8$, and the classical
\textbf{bootstrap exceedance share} --- the fraction of resamples in which $\hat\tau^*(x)$
happened to exceed \eur{\nbOneCost} --- thresholded at the identical $0.8$. They finish
\emph{level}. The posterior rule earns \nbOneBayesEdge{} euros against the best classical rule on
the board: \pc{\nbOneBayesEdgePct} of the oracle prize.

One objection has to be met before that number is banked, because the ladder has just been unkind to
the model it is read off. The euros above come from BCF's posterior, and BCF is the arm the ladder
debits --- including on top-decile coverage, which is the very quantity $\Prob(\tau_i > c)$
integrates. So \cref{tab:up:rules} prices the identical rule on the \textbf{BART T-learner's}
posterior as well: it contacts \nbOneBartRuleN{} customers, realises \eur{\nbOneBartRuleProfit}, and
beats the best classical rule by \nbOneBartRuleEdge{} euros. The Bayesian arm's \emph{best} euro
number, not merely its first, is therefore on the board --- and the conclusion is unmoved. BCF
remains the chapter's decision model because its propensity input is what targeted selection would
require, and because nothing in the euros turns on the choice.

\textbf{On this decision, at this $n$, the Bayesian layer bought essentially no extra euros. Say
it out loud.} Every confidence-aware rule also beats targeting on either arm's bare \emph{point}
estimate, which over-contacts (\nbOnePointRuleN{} customers against the honest rule's
\nbOneBayesRuleN) because a point estimate cannot tell a confident \eur{9} from a coin-flip
\eur{9}. That gap --- between point-targeting and confidence-targeting --- is real and worth euros.
The gap between the \emph{Bayesian} and the \emph{classical} way of being confidence-aware is not.

\subsection*{The durable reason to prefer the posterior: a probability you can act on}

It is \cref{tab:up:calibration}, and it is not a philosopher's quibble.

\input{build/tables/nb01_calibration}

Both arms hand the decision rule a number in $[0,1]$ per customer. Bucket customers by it and ask
what share of them truly had $\tau_i > c$. Neither number is perfectly calibrated --- both are
optimistic in the middle buckets --- but only one of them is \emph{trying} to be a probability.

$\Prob(\tau_i > c \mid D)$ is a statement about \emph{this customer's effect}, obtained by
integrating a posterior, and it is the object that \cref{sec:up:euro}'s rule, its confidence
sweep, its straddler set and its value-of-information integral are every one of them
\emph{defined as an expectation over}. The bootstrap exceedance share is a statement about the
\emph{estimator}: how often a refitted T-learner would land above \eur{\nbOneCost} if the customer
list were redrawn. It is centred on $\hat\tau(x)$, not on $\tau_i$; it inherits the learner's
shrinkage in every single replicate, so it structurally cannot warn anyone that it is biased ---
the top-decile coverage above is that failure, made visible; and nothing licenses reading it as a
degree of belief about a person. \Cref{ch:pvp} makes the general statement; this table is what
it looks like when a marketing team runs into it.

That it \emph{performs} nearly as well here as a decision heuristic is the honest and
uncomfortable finding, and this chapter reports it rather than hiding it. What it cannot do is be
\emph{interpreted}, \emph{dialled} against a stated risk appetite, or fed to a value-of-information
integral without silently changing the meaning of every euro that comes out.

\subsection*{So which model does the rule trust?}

Assemble the findings as a decision rather than a scorecard. The rule ships on
$\Prob(\tau_i > c \mid D)$, a per-customer probability, so the chosen model has to do two things:
rank customers well, and produce that probability honestly. On the ranking, \cref{tab:up:ladder} is
decisive --- the accuracy is real and it belongs to the \emph{base learner}, a sum-of-trees
posterior in place of a point fit, while BCF's specifically causal package adds no ranking accuracy
and, as the euros above showed, no euros. On the probability, \cref{tab:up:calibration} is decisive
the other way --- only a posterior yields $\Prob(\tau_i > c \mid D)$ at all, and the classical
exceedance share, however close its decisions land here, is simply not that object. So the rule runs
on a Bayesian posterior, and it runs there for the base learner and the decision quantity, not for
the two ingredients still on probation: with the confounders fully observed the propensity input has
nothing to fix, and the shrink-toward-homogeneity prior costs the top decile its coverage in exactly
the place the rule spends money. Adjustment, not Bayes, is what defeated the confounding
(\cref{sec:up:classical}); a hidden confounder would drag the posterior exactly as far as it drags
OLS (\cref{sec:up:wrong}); and on a smaller list or a noisier outcome the base learner's averaging
margin only grows. \Cref{sec:up:euro} prices all of this in euros and turns it into a test.

% =====================================================================================
\section{The decision, in euros}
\label{sec:up:euro}

\subsection*{The honest targeting rule}

Targeting on ``posterior mean above cost'' is a mistake with a euro price. A customer whose point
estimate says \eur{9} but whose interval straddles the \eur{\nbOneCost} cost line is a coin flip,
not a yes. So the rule this chapter recommends thresholds the \emph{probability}, not the estimate:
%
\begin{equation}
  \pi(x_i) \;=\; \mathbf 1\Big[\ \Prob\big(\tau_i > c \;\big|\; D\big) \;>\; \nbOneConfBar\ \Big],
  \qquad
  \Prob\big(\tau_i > c \mid D\big) \;\approx\; \frac{1}{S}\sum_{s=1}^{S}
      \mathbf 1\big[\,\tau_i^{(s)} > c\,\big],
  \label{eq:up:rule}
\end{equation}
%
where $\tau_i^{(s)}$ are the posterior draws of customer $i$'s effect. The threshold $0.8$ is a
convention and can be argued with --- it is a statement of risk appetite, and
\cref{fig:up:voi} prices moving it. The \emph{object} on the left cannot be argued with: it is
\Cref{ch:pvp}'s $\Prob(\tau > c \mid D)$, per customer, and the classical arm
cannot produce it.

Applied here, \cref{eq:up:rule} contacts \nbOneNTargeted{} of \nbOneNCustomers{} customers
(\pc{\nbOneTargetShare}) and realises \eur{\nbOneModelProfit} --- \pc{\nbOneCapture} of the
oracle's \eur{\nbOneOracleProfit}, against the \eur{\nbOneTreatAllLoss} that emailing everyone
\emph{loses} (\cref{fig:up:policy}).

\input{build/floats/nb01_policy}

\subsection*{The profit curve, and where to stop}

Rank customers by predicted effect, contact down the list, and plot cumulative realised
$(\tau_i - c)$. It climbs while persuadables are being added and falls once the list runs into
sure things and sleeping dogs, so \textbf{its peak is the optimal stopping point}: contact the top
\pc{\nbOneStopFrac}, earn \eur{\nbOneStopProfit}. This is the single most useful artefact in the
chapter for an operator, because it converts a model into an instruction --- ``email the top
\pc{\nbOneStopFrac} of the ranked list'' --- and prices what the next customer down is worth.

Three euro figures in this section sit within a few euros of one another, and they are \emph{three
different objects}. Do not read them as one number quoted thrice. The shipped rule of \cref{eq:up:rule}
thresholds a \emph{probability} at \nbOneConfBar{} and realises \eur{\nbOneModelProfit} --- a policy
you can run on the real list, because $\Prob(\tau_i > c \mid D)$ survives production. This curve instead
ranks by the \emph{point} prediction and reads its peak off the \emph{known} truth, so it lands a
little higher, at \eur{\nbOneStopProfit}, and is a simulator-only ceiling, not a shippable rule. The
confidence sweep below traces the shipped rule's value as the bar turns, peaking on a plateau at that
same \eur{\nbOneModelProfit}. Rule value, hindsight ceiling, sweep plateau: near in euros, distinct in
kind.

Note what the curve is \emph{not}: it is not a forecast. Its shape is realised on the known truth,
which is available here and nowhere else. In production the same curve is estimated from a
randomized holdout by the Qini construction of \cref{sec:up:valid}, and it is noisier, especially
at the top where the slices are thin.

\subsection*{What the remaining uncertainty is worth --- and to whom}

\textsc{Test further} must not end as a shrug. The deliverable is an answer to \emph{what evidence
would settle this, and is it worth buying?} Uncertainty has a price: the euros lost by acting on a
noisy estimate rather than the truth. Per customer,
%
\begin{equation}
  \text{VOI}_i \;=\; \E\big[\max(\tau_i - c,\ 0)\big] \;-\; \max\big(\E[\tau_i] - c,\ 0\big) ,
  \label{eq:up:voi}
\end{equation}
%
the gain from deciding \emph{after} the uncertainty resolves rather than committing now on the
posterior mean \citep{raiffa1961}. \Cref{eq:up:voi} is zero for anyone whose posterior sits
confidently on one side of the cost line --- if the answer will not change, the information is
worthless --- and it is large exactly for the \textbf{straddlers}, whose credible interval crosses
$c$.

And that is what the notebook finds: of \eur{\nbOneVoiTotal} of total value of information,
\eur{\nbOneVoiStraddlers} sits on the \nbOneNStraddlers{} straddlers --- \pc{\nbOneStraddlerShare}
of the list (\cref{fig:up:voi}). \textbf{That is not an abstraction; it is a test design.} The
follow-up is a small randomized A/B test \emph{on the straddlers}, not a blanket send, and
\cref{eq:up:voi} bounds what such a test can be worth \emph{under this posterior} --- it prices
\emph{perfect} information, and a finite experiment resolves only part of the uncertainty.

That scope clause is doing work, and it cuts the other way from the one a reader expects. Both the
integral and the straddler set are read off BCF's own draws, and \cref{sec:up:compare} measured those
draws as under-covering the top decile of true effect at \pc{\nbOneBcfTopDec} against a nominal
90\,\%. Intervals that are too narrow cross the cost line \emph{less} often than honest ones would,
so \nbOneNStraddlers{} straddlers and \eur{\nbOneVoiStraddlers} of value are a \textbf{floor}: a
better-calibrated posterior would report more genuinely undecided customers, not fewer, and the test
would be worth more, not less. The number is safe to spend against; it is not safe to economise
against. What it does rule out is the other error: a test costing many times
\eur{\nbOneVoiStraddlers} to settle an \eur{\nbOneVoiStraddlers} question is not caution. It is an
expensive way to feel careful.

\input{build/floats/nb01_voi}

\subsection*{The price of insisting on certainty}

The confidence bar in \cref{eq:up:rule} is a dial, and \cref{fig:up:voi} (left) shows what turning
it costs. Set it too low and the rule buys coin flips: at a bar of \nbOneConfThrLoose{} it realises
\eur{\nbOneConfProfitLoose}. Set it too high and it refuses profitable customers out of
squeamishness: at \nbOneConfThrStrict{} it realises \eur{\nbOneConfProfitStrict}. The chapter's
$0.8$ sits on the plateau between them, at \eur{\nbOneModelProfit}, and \emph{that} is the honest
reason to prefer it --- not because $0.8$ is principled, but because the answer is insensitive to
it across a wide range. Where a decision \emph{is} sensitive to an arbitrary knob, that
sensitivity is the finding and must be reported as one.

\begin{remark}[The uncomfortable inheritance]
\Cref{eq:up:rule} is read off the same posterior whose top-decile coverage
\cref{sec:up:compare} just measured at \pc{\nbOneBcfTopDec}. The decision quantity lives in the arm
whose uncertainty is least trustworthy \emph{exactly where the rule spends money} --- the same
tension \Cref{ch:geo} meets on a geo panel and states as a rule: use the posterior for the
decision, and design-based inference for your confidence in it. Here that means: ship the
targeting rule, and \emph{test the straddlers}. The straddler experiment is not a nicety. It is
the design-based leg of a decision whose model-based leg is known to be overconfident.
\end{remark}

% =====================================================================================
\section{What can go wrong}
\label{sec:up:wrong}

\subsection*{The S-learner, and the discipline of failing precisely}

\Cref{tab:up:bakeoff} runs three meta-learners --- S, T and BCF, one BART base learner throughout
--- on both the randomized twin and the confounded panel.

\input{build/tables/nb01_bakeoff}

The S-learner is the clear loser, and \emph{what} it loses is the lesson. Its single shared surface
regularises the treatment column toward zero, so it flattens $\tau(x)$ even with no confounding to
fight (\cref{fig:up:flatten}, left). Against the BART T-learner --- same base learner, so the
comparison isolates the ``one shared surface'' choice and nothing else --- the verdicts are:
\textbf{\nbOneSflatPeheWin} on PEHE (\nbOneSflatPeheA{} against \nbOneSflatPeheB),
\textbf{\nbOneSflatAuucWin} on AUUC (\nbOneSflatAuucA{} against \nbOneSflatAuucB), and
\textbf{\nbOneSflatCovWin} on coverage --- where the S-learner's intervals cover
\pc{\nbOneSlCov} against a nominal 90.

\input{build/floats/nb01_flatten}

Three things follow, and they are all counter-intuitive.

\textbf{(1) The ranking survives.} The S-learner still captures \pc{\nbOneSlAuucPct} of the
oracle's ranking value, and with the GBM base learner of \cref{sec:up:classical} it cost the
ranking essentially nothing at all (\cref{tab:up:cate}). A uniformly squashed $\tau(x)$ can still
put customers in nearly the right order. \textbf{The severity of the S-learner's failure is a
property of the base learner, not of the letter S} --- which is why this chapter reports it axis
by axis instead of repeating the slogan.

\textbf{(2) The intervals do not survive.} \pc{\nbOneSlCov} coverage against a nominal 90 is not a
dent; it is a collapse. Flattening always destroys magnitudes, sometimes dents the ranking, and
destroys honesty every time.

\textbf{(3) The average survives, and that is the trap.} The S-learner's ATE bias is
\nbOneSlAteBias{} euros, against the naive difference in means' \nbOneNaiveAteBias{} on the same
panel --- an order of magnitude closer. \textbf{An ATE-only dashboard would have shipped it.} The
S-learner's fault is heterogeneity, not the average, and a monitoring stack that watches only the
average will never see it.

\subsection*{The assumption that cannot be tested}

\Cref{as:up:unconf} is untestable, so the honest question is not ``is it true?'' but ``how fragile
is the conclusion to a confounder that was missed?'' \Cref{fig:up:sensitivity} answers it two ways.

\textbf{The tipping point.} A hidden driver $U$ is planted into both the emailing rule and spend,
its strength swept, and the effect re-estimated adjusting for the \emph{observed} features only,
blind to $U$ --- exactly as a real analyst would be. The estimate crosses the \eur{\nbOneCost}
cost line at a strength of about \nbOneTipping. Beyond that, confounding \emph{alone} would fake
an ``email everyone pays'' signal out of nothing. The two-dimensional version, genuinely refit at
every grid point, shows the geometry: omitted-variable bias needs \emph{both} channels, scaling
with the product of $U$'s pull on treatment and its pull on spend, so the frontier is a hyperbola
and not a slope.

\textbf{The E-value.} \citet{vanderweele2017}'s E-value is the minimum strength of association, on
the risk-ratio scale, that a hidden confounder would need with \emph{both} treatment and outcome
to explain the estimated effect away entirely. Computed on the AIPW estimate --- never on the
simulation truth, which a real analyst does not have --- and standardized by the outcome standard
deviation, since the E-value is defined multiplicatively and a euro effect must be converted
first, it is \nbOneEvalue{} (\nbOneEvalueLo{} at the near confidence limit). That is
\textbf{moderate, not decisive}: robustness that is real but does not license complacency. The
right response is not more statistics; it is to ask a domain expert whether anything that strong
went unmeasured.

\input{build/floats/nb01_sensitivity}

\subsection*{And the ones no diagnostic here would catch}

\textbf{Overlap can fail silently in production} even though it holds here. A segment that was
\emph{always} emailed produces a propensity pinned at 1, no comparison group, and an extrapolation
dressed as an estimate. Trim it or declare it.

\textbf{SUTVA} (\cref{as:up:sutva}) breaks the moment a discount code leaks. Nothing in this
chapter would notice.

\textbf{And the estimator itself is not reproducible.} \code{pymc-bart} is not seed-stable in this
environment, which is why no BART-derived number appears as a literal anywhere above. That is
uncomfortable and it is stated rather than smoothed over: the \emph{decisions} in
\cref{tab:up:rules} are robust across runs, and several of the \emph{metric comparisons} are not,
which is exactly why the latter are reported with a tolerance and are permitted to say ``too close
to call''.

% =====================================================================================
\section{Takeaways}
\label{sec:up:take}

\begin{enumerate}[leftmargin=*, itemsep=.45em]
  \item \textbf{The ATE is the wrong object, whichever way it points.} The average effect is
        \eur{\nbOneTrueAte}, correctly estimated by three classical estimators, and it issues only two
        instructions. Below the cost it says email nobody, forgoing the whole \eur{\nbOneOracleProfit}
        the ranking would earn. Above the cost --- what the naive \eur{\nbOneNaiveEst} says --- it says
        email everyone, which earns $n(\text{ATE} - c)$, here \eur{\nbOneTreatAllLoss} \emph{destroyed}.
        A better average repairs neither, because \pc{\nbOneShareProfitable} of the list clears the cost
        while \pc{\nbOneShareSleeping} are sleeping dogs the campaign harms, and the two want opposite
        treatment. The business object is the CATE $\tau(x)$ and the \emph{ranking} it induces; the
        deliverable is a policy: mail the top-$k$ (\cref{tab:up:rules}).

  \item \textbf{Heterogeneous effects do not require Bayes.} A T-learner --- two regressions and a
        subtraction --- ranks customers at AUUC \nbOneGbmTAuuc{} of the oracle, with a bootstrap band
        for free. It needs the identification argument (\crefrange{as:up:unconf}{as:up:sutva}) and a
        base learner flexible enough for the true $\tau(x)$: swap the trees for a straight line and
        AUUC collapses to \nbOneOlsTAuuc.

  \item \textbf{You cannot score a CATE model directly, and Qini is why uplift works in production.}
        $\tau_i$ is never observed, so PEHE is a luxury of the simulator. AUUC/Qini borrows the
        randomization: within the top-$k$ slice, treated and control are exchangeable, so their
        difference in means estimates that slice's effect with no counterfactual. It grades the
        \emph{ranking} only, and is blind to a model that flattens every magnitude
        (\cref{fig:up:flatten}).

  \item \textbf{Attribute the win before you bank it.} The ladder (\cref{tab:up:ladder}) prices one
        rung at a time, and its verdict is that the Bayesian edge over the classical arm is
        \nbOneAttrVerdict. In numbers: the base learner buys \nbOneAttrBaselearner; BCF's package ---
        its shrink prior, its propensity input and its joint additive surface, which move together and
        are judged together --- buys \nbOneAttrBcfBuys{} and \emph{costs} \nbOneAttrBcfCosts. With the
        confounders fully observed and overlap good, the propensity trick has nothing to fix.
        \textbf{Do not credit a prior for what the base learner produced}, and do not credit one
        ingredient for the verdict on a bundle of four.

  \item \textbf{The intervals are worst exactly where the money is.} BCF's 90\,\% credible intervals
        cover \pc{\nbOneBcfCovNinety} of customers overall but only \pc{\nbOneBcfTopDec} on the top
        decile of true effect, against the bootstrap's \pc{\nbOneBootTopDec}; the head-to-head names
        \nbOneHhTopdecWin. Shrinkage pulls the biggest effects down, and the interval --- of ordinary
        width (\eur{\nbOneBcfTopdecWidth}, \nbOneTopdecWidthVerdict{} the classical band's
        \eur{\nbOneBootTopdecWidth}) --- is centred on the shrunken value and does not reach the truth.
        The model is most confident, and most wrong, about the customers it is about to spend money on.

  \item \textbf{On the euros, the Bayesian layer bought almost nothing --- and the reason to prefer it
        is not the euros.} The posterior rule beats the best classical rule by \nbOneBayesEdge{} euros,
        \pc{\nbOneBayesEdgePct} of the oracle prize; the same rule on the BART T-learner's posterior ---
        the model the ladder prefers --- beats it by \nbOneBartRuleEdge{} euros (\cref{tab:up:rules}),
        so the verdict does not turn on which Bayesian model is asked. What the posterior alone
        produces is $\Prob(\tau_i > c \mid D)$: a probability about \emph{the customer}, which can be
        thresholded against a risk appetite, swept, and fed to a value of information. The bootstrap
        exceedance share performs nearly as well and \emph{is not that object}
        (\cref{tab:up:calibration}); it is the scatter of an estimator across datasets never collected,
        and it inherits its own shrinkage in every replicate.

  \item \textbf{Act where you are sure; measure where you are not.} \eur{\nbOneVoiStraddlers} of the
        \eur{\nbOneVoiTotal} value of information sits on the \nbOneNStraddlers{} straddlers whose
        interval crosses the cost line (\pc{\nbOneStraddlerShare} of the list). So the follow-up is a
        randomized test \emph{on the straddlers}. That test is also the design-based leg the
        overconfident top decile demands, and, because that overconfidence narrows the intervals that
        define a straddler, a \emph{floor} on the test's worth, not a ceiling. Ship the rule; test the
        doubt.
\end{enumerate}

% =====================================================================================
\clearpage
\section*{Appendix 3.A\quad Formal machinery}
\label{app:up:tol}
\addcontentsline{toc}{section}{Appendix 3.A\quad Formal machinery}

The body states each idea plainly, gives the euro result, and points here. This appendix carries the
formal version of each --- the estimand ladder, the meta-learners, BCF's parameterisation, the grading
metrics, and the tree-sampler tolerances --- every one with the worked example the body refers to.
Nothing here is needed to follow the chapter's decisions. (The simulator constants live in
\cref{tab:up:params}.)

\subsection*{3.A.1\quad The CATE and the estimand ladder}

Write $Y_i(1)$ for customer $i$'s spend if emailed and $Y_i(0)$ for their spend if not; the individual
effect is $\tau_i = Y_i(1) - Y_i(0)$, and only one of its two terms is ever observed. Three averages
are built on it, and the chapter needs all three:
%
\begin{align}
  \text{ATE} &\;=\; \E\big[Y(1) - Y(0)\big]
     &&\text{the effect on a random customer,} \label{eq:up:appate}\\
  \tau(x) &\;=\; \E\big[Y(1) - Y(0) \,\big|\, X = x\big]
     &&\text{the CATE: the effect for customers like } x, \label{eq:up:appcate}\\
  V(\pi) &\;=\; \E\big[(\tau(X) - c)\,\pi(X)\big]
     &&\text{the euros a targeting rule } \pi \text{ earns.} \label{eq:up:apppolicy}
\end{align}
%
The ATE prices \emph{roll it out to everyone?}; the CATE decides \emph{who next?}; the policy value
$V(\pi)$ is what the business buys, and it rewards ranking, not magnitude. The rungs are not
interchangeable. Here the ATE is \eur{\nbOneTrueAte}, below the \eur{\nbOneCost} cost, so an
average-level rule says email no one --- yet the CATE runs from \eur{\nbOneTauLo} to \eur{\nbOneTauHi}
across customers, and a rule that mails only where $\tau(x) > c$ earns the \eur{\nbOneOracleProfit} the
average leaves on the table.

\subsection*{3.A.2\quad The meta-learners: S, T, and X}

A meta-learner turns any regression routine into a CATE estimator by how it arranges the fits. All
three below estimate the same $\tau(x)$ of \cref{eq:up:appcate}; they differ only in construction.

\begin{description}[leftmargin=*, style=nextline]
  \item[T-learner.] Fit one regression per arm and subtract,
        $\hat\tau(x) = \hat\mu_1(x) - \hat\mu_0(x)$, with $\hat\mu_t$ fitted on the customers with
        $T = t$. Two independent fits.
  \item[S-learner.] Fit \emph{one} regression on $(x, T)$ jointly and read
        $\hat\tau(x) = \hat f(x, 1) - \hat f(x, 0)$. Treatment is one more feature, so a regulariser
        can shrink its influence and flatten $\tau(x)$.
  \item[X-learner.] Fit the two arms, impute each customer's missing counterfactual from the other
        arm's fit to form imputed effects $\tilde\tau_i$, regress those on $x$ within each arm, and
        blend the two by the propensity. Built to help when the arms are very unequal in size.
\end{description}
%
The construction, not the letter, decides what breaks. With the GBM base learner the T-learner reaches
AUUC \nbOneGbmTAuuc{} of the oracle; the S-learner's shared surface flattens the effect, so its PEHE is
\eur{\nbOneGbmSPehe} against the T-learner's \eur{\nbOneGbmTPehe}, while its ranking barely suffers.

\subsection*{3.A.3\quad BCF: the treatment-surface parameterisation}

A Bayesian Causal Forest \citep{hahn2020} splits the mean into a \emph{prognostic} surface and a
\emph{treatment} surface, each a sum-of-trees BART prior, and reads the CATE straight off the second:
%
\begin{equation}
  \mu_i \;=\; f_{\text{prog}}\big(x_i,\ \hat e(x_i)\big) \;+\; f_{\tau}(x_i)\, T_i,
  \qquad
  \hat\tau(x) \;=\; f_\tau(x),
  \label{eq:up:appbcf}
\end{equation}
%
with the estimated propensity $\hat e(x)$ fed to $f_{\text{prog}}$ to absorb targeted selection, and
$f_\tau$ given \emph{half} as many trees as $f_{\text{prog}}$ so heterogeneity is reported only when
the data insist. \Cref{eq:up:bcfprior} is the full prior. Here $f_{\text{prog}}$ carries
$m = \nbOneBcfTrees$ trees and $f_\tau$ carries $m/2$; the posterior-mean ATE is \eur{\nbOneBcfAte},
agreeing with the classical arm.

\subsection*{3.A.4\quad The Qini / uplift curve and AUUC}

Sort customers by \emph{predicted} effect, contact them from the top down, and plot the cumulative
\emph{true} incremental effect against the fraction contacted. Random targeting traces the diagonal;
an oracle that knows every $\tau_i$ traces the upper envelope. The \textbf{AUUC} is the area between
the model's curve and the diagonal, normalised by the oracle's:
%
\begin{equation}
  \text{AUUC}
  \;=\; \frac{\int_0^1 \big(Q_{\text{model}}(p) - p\,Q(1)\big)\,dp}
             {\int_0^1 \big(Q_{\text{oracle}}(p) - p\,Q(1)\big)\,dp},
  \label{eq:up:appqini}
\end{equation}
%
where $Q(p)$ is the cumulative true effect captured by contacting the top fraction $p$. Read it as a
share: $1.0$ is a perfect ranking, $0$ is no better than random. It scores only the \emph{order},
\cref{eq:up:apppolicy}'s sole input, and is blind to magnitude. BCF ranks at AUUC \nbOneBcfAuuc{} of
the oracle's (\cref{fig:up:qini}).

\subsection*{3.A.5\quad PEHE: the CATE mean-squared error}

PEHE --- Precision in Estimation of Heterogeneous Effects --- is the root-mean-square error of the
estimated effect against the true one:
%
\begin{equation}
  \text{PEHE} \;=\; \sqrt{\frac1n \sum_{i=1}^{n}\big(\hat\tau(x_i) - \tau_i\big)^2},
  \label{eq:up:apppehe}
\end{equation}
%
in euros; lower is better. It is the CATE's answer to RMSE, and because it compares against $\tau_i$ it
exists \emph{only in a simulator} --- \cref{sec:up:valid} says what replaces it in production. Here BCF
scores \eur{\nbOneBcfPehe} against effects that run from \eur{\nbOneTauLo} to \eur{\nbOneTauHi}, and its
estimates correlate with the truth at \nbOneBcfCorr.

\subsection*{3.A.6\quad Calibration and coverage of the CATE intervals}

An interval is \textbf{calibrated} at level $1 - \alpha$ when the true effect lands inside it that
share of the time. \textbf{Coverage} measures it:
%
\begin{equation}
  \text{coverage}
  \;=\; \frac1n \sum_{i=1}^{n} \mathbf 1\big[\,\tau_i \in \text{CI}_{1-\alpha}(x_i)\,\big],
  \label{eq:up:appcov}
\end{equation}
%
reported overall and by decile of the true effect, because a model can be calibrated on average and
overconfident exactly where the money is. Here BCF's 90\,\% intervals cover \pc{\nbOneBcfCovNinety}
overall but only \pc{\nbOneBcfTopDec} on the top decile (\cref{tab:up:calibration}) --- a failure of
location, not width, as shrinkage pulls the largest effects down.

The remaining machinery is specific to the tree sampler. It exists because \code{pymc-bart} is not
seed-stable in this environment: two full runs of identical source have moved this model's PEHE by more
than a fifth and flipped a coverage-winner line. So no BART-derived number is typed into the chapter
--- every one is injected from the run that produced this PDF --- and no two are compared by eye.

\subsection*{3.A.7\quad The convergence verdict for a tree sampler}

Divergences are a NUTS concept and do not apply to the particle-Gibbs sampler PGBART; what remain
are $\hat R$ and the effective sample size. Against those this book applies a fixed band to a tree
sampler: $\hat R \le 1.01$ is a clean pass, $1.01$--$1.02$ is tolerated, and anything above that is
an \emph{investigate} signal. The verdict is computed and printed rather than eyeballed so it cannot
drift with the author's optimism.

The band has to be read with a caveat that is specific to sum-of-trees models. The scalar being
monitored is the observation-noise standard deviation $\sigma$, and it is coupled to a sum of
\nbOneBcfTrees{} trees whose \emph{individual} structures are exchangeable. A chain can therefore
explore a perfectly good posterior \emph{over functions} while its scalar summaries mix slowly, so
raising the draw count grinds $\hat R$ down only slowly, and $\hat R$ is a poor proxy for the object
the chapter actually reads off the model --- the treatment surface $f_\tau$. The honest response is
not to wave the number through but to reach for a different backstop: the \nbOneNSeeds-seed refit of
\cref{sec:up:valid}, which grades the \emph{estimates} rather than the sampler and finds them
\nbOneStabVerdict.

\paragraph{Worked example.} On the run behind this PDF the headline BCF fit returns
$\hat R = \nbOneBcfRhat$, and the band above turns that into the printed verdict
\textbf{\nbOneBcfRhatVerdict}. The chapter then does exactly what a cautious reading of a tree
sampler prescribes: it does not trust the scalar on its own, it leans on the multi-seed refit, and
it treats the treatment surface as the thing to grade. A converged sampler would not have rescued
the deeper problem in any case, because convergence is not reproducibility: the seed-instability
above would remain.

\subsection*{3.A.8\quad The comparison tolerance}

Every head-to-head between two BART-derived numbers is judged against a tolerance, and a gap is
called a winner only when it exceeds that tolerance; otherwise the verdict is allowed to read
``too close to call''. A comparison rule that \emph{cannot} return that answer will always
manufacture a winner out of noise. The tolerance sums two sources of noise.

\begin{description}[leftmargin=*, style=nextline]
  \item[The paired bootstrap.] Resample the same customers for \emph{both} arms, refit each, and
        watch the gap between their scores move; the base tolerance is twice the standard error of
        that gap. Pairing on the customer list is what makes it a fair test of the \emph{difference}
        rather than of two independent estimates.
  \item[The sampler wobble.] The paired bootstrap resamples data; it cannot see the noise the
        \code{pymc-bart} sampler injects at fixed data. So wherever a BART \emph{refit} sits on one
        side of a comparison, the fit-to-fit wobble is added on top. Its size is measured directly:
        two equal-compute refits of the identical specification, same customers and a different seed,
        disagree by \eur{\nbOneBcfSeedPeheWobble} of PEHE and \nbOneBcfSeedCovWobble{} points of
        overall coverage, and that disagreement is pure sampler noise by construction. (It is also
        why \cref{tab:up:compare} and \cref{tab:up:ladder} differ: they are two fits of one model.)
        \Cref{fig:up:seedwobble} plots that floor across the multi-seed refits.
\end{description}

\input{build/floats/nb01_seed_wobble}

\paragraph{Worked example.} Take the headline PEHE head-to-head of \cref{sec:up:compare}: BCF at
\nbOneHhPeheA{} against the GBM T-learner at \nbOneHhPeheB, a gap of \nbOneHhPeheGap. The tolerance
assembled from the two sources above is \nbOneHhPeheTol. The verdict is whatever the comparison of
those two quantities returns --- here \textbf{\nbOneHhPeheWin} --- computed by the notebook and
cited, never argued. Read every ``\dots against a tolerance of \dots'' clause in
\cref{sec:up:compare} as this same construction applied once more.

\notebooknote{%
  This chapter is \code{notebooks/01\_uplift\_targeting.ipynb}. The simulator is
  \code{cmp.dgp.uplift\_customers} (\crefrange{eq:up:mu}{eq:up:y}); the classical arm is
  \code{cmp.classical} plus scikit-learn's gradient boosting; the Bayesian arm is
  \code{cmp.estimators.bcf}, a \code{pymc-bart} model. The notebook runs in a fast teaching mode
  (\code{CMP\_FAST=1}) and a full mode (\code{CMP\_FAST=0}); every number in this chapter comes from
  the full run, emitted by \code{cmp.report} and injected here as a macro --- which in this chapter
  is a load-bearing property rather than a tidiness one, because \code{pymc-bart} is not
  seed-reproducible in this environment and any BART number typed into prose would be false by the
  next execution. The notebook also carries a network-gated section (\code{CMP\_REAL=1}) that runs
  the identical pipeline on Hillstrom's public 64{,}000-customer randomized email experiment
  \citep{hillstrom2008}, where the ATE is trustworthy without adjustment, no per-customer truth
  exists, and the Qini curve of \cref{sec:up:valid} is therefore the only grade available --- which
  is exactly what production looks like.}
